\section{Details of the explicit lattice Green bound}
\label{app:green-kernel}

We justify the Fourier calculation in Lemma~\ref{lem:green-kernel} with
the canonical bond-anchor convention.  Orthogonality of a field
$f\in\mathcal H$ to $Dv$ for every compact $v$ is equivalent, after
Fourier transform, to
$\sum_i\overline{d_i}\widehat f_i=0$ almost everywhere: trigonometric
polynomials are dense and the $d_i$ are bounded.  Thus $D(\X)$ is the
pointwise span of $d=(d_1,d_2)$, apart from the irrelevant frequency zero,
and the projection symbol is
\[
 m_{ij}(k)=\frac{d_i(k)\overline{d_j(k)}}{\Lambda(k)},\qquad
 \Lambda(k)=4-2\cos k_1-2\cos k_2.
\]
This argument does not require an $\ell^2$ vertex potential for each
finite-energy gradient.  In particular, $\nrm{\mathsf P}_{\mathcal H\to
\mathcal H}=1$ and every scalar symbol is bounded by one.

Let $z=(n_1,n_2)$ and first assume
$r=|n_2|=\nrm{z}_\infty\ge1$.  For $t=k_1\ne0$, set
\[
 s=|\sin(t/2)|,\quad
 \alpha=2\operatorname{arsinh}s,\quad
 g_n(t)=\frac{e^{-\alpha|n|}}{2\sinh\alpha}.
\]
Since $\cosh\alpha=2-\cos t$, the Poisson kernel identity gives
\begin{equation}
 \frac1{2\pi}\int_{-\pi}^{\pi}
 \frac{e^{\mathrm i n y}}{4-2\cos t-2\cos y}\,dy=g_n(t).
 \label{eq:green-partial}
\end{equation}
It follows by expanding
$(\cosh\alpha-\cos y)^{-1}$ as a geometric Fourier series; all series
converge absolutely for this fixed $t\ne0$.

The partially integrated diagonal entries are
$4s^2g_{n_2}$ and $\delta_{n_2,0}-4s^2g_{n_2}$.  Since $n_2\ne0$,
their magnitudes are at most
\[
 4s^2g_{n_2}
 =\frac{s}{\sqrt{1+s^2}}e^{-\alpha r}
 \le s e^{-\alpha r}.
\]
For the $12$ entry, the partial coefficient is
$d_1[g_{n_2-1}-g_{n_2}]$; for the $21$ entry it is
$\overline{d_1}[g_{n_2+1}-g_{n_2}]$, up to the choice of Fourier sign.
For both shifts, including $r=1$,
\begin{equation}
 |g_{n_2\pm1}-g_{n_2}|
 \le\frac{e^\alpha-1}{2\sinh\alpha}e^{-\alpha r}
 =\frac{e^\alpha}{e^\alpha+1}e^{-\alpha r}
 \le e^{-\alpha r}.
 \label{eq:green-shift-bound}
\end{equation}
As $|d_1|=2s$, the two mixed entries have partial magnitude at most
$2s e^{-\alpha r}$.  Thus no unit shift of the anchor distance remains
unaccounted for.

For $0\le t\le\pi$,
$s\le t/2$, $s\ge t/\pi$, and
$\operatorname{arsinh}s\ge s/\sqrt2$.
The last inequality follows by integrating
$(1+s^2)^{-1/2}\ge1/\sqrt2$ on $[0,s]\subset[0,1]$.
Consequently
\[
 \frac1{2\pi}\int_{-\pi}^{\pi}s e^{-\alpha r}\,dt
 \le\frac1{2\pi}\int_0^\infty
      t e^{-\sqrt2 rt/\pi}\,dt
 =\frac{\pi}{4r^2}.
\]
This gives $|\mathsf P_{ij}(z)|\le\pi/(2r^2)<2/r^2$.
If $|n_1|=\nrm{z}_\infty\ge1$, exchange the axes and their orientation
labels.  At $z=0$ use the integral of the bounded scalar symbol.
The multiplier entries are bounded on a finite measure torus, so Fubini
applies directly.  The exceptional value $t=0$ has measure zero, and the
partial bounds above are integrable.  There is no appeal to an
uncontrolled logarithmically divergent vertex Green function.

\section{An annular Poincar\'e constant for the cutoff}
\label{app:green-annulus}

For $R\ge1$ integer and
$\mathcal A_R=\{x:R\le\nrm{x}_\infty\le2R\}$, we prove
\begin{equation}
 \sum_{\mathcal A_R}|u-\overline u_{\mathcal A_R}|^2
 \le\frac{225}{4}R^2\sum_{e\subset\mathcal A_R}|D_eu|^2.
 \label{eq:green-annular-poincare}
\end{equation}
For $n$ values on an interval, the variance identity and Cauchy--Schwarz
give
\[
 \sum_i|u_i-\overline u|^2
 =\frac1n\sum_{i<j}|u_i-u_j|^2
 \le\sum_{k=1}^{n-1}\frac{k(n-k)}2|u_{k+1}-u_k|^2
 \le\frac{n^2}{8}\sum_k|D u_k|^2.
\]
Decomposing the variance of an $n_1\times n_2$ rectangle into row
variances and the variance of row means, followed by Jensen's inequality,
gives the constant $\max(n_1^2,n_2^2)/8$ for its full edge energy.

Cover the annulus by its four top, right, bottom, and left strips, each
with $(4R+1)(R+1)=N$ vertices.  Adjacent strips have an overlap of
$J=(R+1)^2$ vertices.  Write $m_i$ for their means and
$V_i=\sum_{S_i}|u-m_i|^2$.  Averaging over an overlap yields
$|m_i-m_j|^2\le2(V_i+V_j)/J$.  If $c$ is the mean of the four $m_i$,
the four-cycle inequality gives
\[
 \sum_i|m_i-c|^2\le\frac12\sum_{i\sim j}|m_i-m_j|^2
 \le\frac2J\sum_i V_i.
\]
Hence
\[
 \sum_{\mathcal A_R}|u-c|^2
 \le\sum_iV_i+N\sum_i|m_i-c|^2
 \le(1+2N/J)\sum_iV_i\le9\sum_iV_i.
\]
Each edge belongs to at most two strips and $(4R+1)^2\le25R^2$.
The rectangular estimate bounds the last expression by
$9\cdot2\cdot25R^2/8$ times the annular edge energy.  Replacing $c$
by the annular mean can only decrease the variance, proving
\eqref{eq:green-annular-poincare}.

\section{Explicit dipolar and payload norm bounds}\label{app:green-inputs}

Let $L=12$ and let $\phi_{12}$ be the real reference in
\cref{sec:green}. For an index radius $r\ge16$, the distance to every
interpolated pole is at least $r-8\ge r/2$. Integrating the derivative
of the argument with respect to pole position, and then along a unit
bond, gives
\[
 |\phi_{12}(x)|\le2L/r,\qquad
 |D_e\phi_{12}|\le4L/r^2
 \quad\text{if }\min(\nrm{x}_\infty,\nrm{y}_\infty)=r.
\]
These integrations take place away from the pole segment, so they do not
introduce branch jumps. There are at most $16r+4\le20r$ bonds with
minimum endpoint radius $r$. Therefore, for integer $R\ge16$,
\[
 \sum_{\min(\nrm{x}_\infty,\nrm{y}_\infty)\ge R}
 |D_e\phi_{12}|^2
 \le320L^2\sum_{r\ge R}r^{-3}\le480L^2/R^2.
\]
For the cutoff $\zeta_R$ used in \cref{lem:green-projection}, the
annulus has $12R(R+1)\le24R^2$ vertices. Thus direct use of its product
identity gives
\[
 \nrm{(1-\zeta_R)\phi_{12}}_\X^2
 \le2(480L^2/R^2)+\frac4{R^2}
           \sum_{\mathcal A_R}|\phi_{12}|^2
 \le1344L^2/R^2.
\]
Since $\sqrt{1344}<40$, best approximation proves
$\nrm{(I-\Pi_{2R})\phi_{12}}_\X\le40L/R=480/R$.

For completeness, a coarse global bound is
\[
 \nrm{\phi_{12}}_\X^2
 \le4\cdot33^2(4\pi)^2+480\cdot12^2/16^2<10^6.
\]
Here the first term overcounts every bond touching $\Lambda_{16}$;
each difference of the two real arguments is at most $4\pi$ in
magnitude. For a compact correction on an $n\times n$ vertex array,
\[
 \nrm{h}_\X\le\sqrt8\nrm h_{\ell^2}\le3n\max|h|.
\]
Exact binary-payload inspection, including finiteness, zero boundary
values, and manifest hashes, gives
\begin{align*}
 \max|h_m|&=\frac{99633080638867}{562949953421312}<.177,\\
 \max|h_s|&=\frac{1834512834580789}{4503599627370496}<.4074,\\
 \max_j\max|z_j|&=\frac{6674417597839975}{2251799813685248}<3.
\end{align*}
The minimum and saddle arrays have side $2048$ and $4096$, respectively;
the stored path displacements $z_j$ have side $256$ and are added to the
minimum candidate. Hence
\begin{align*}
 \nrm{m_0}_\X&<1000+3(2048)(.177)+.017422<6007,\\
 \nrm{s_+}_\X&<1000+3(4096)(.4074)+.012825<6007,\\
 \max_j\nrm{\phi_{12}+h_m+z_j}_\X
   &<1000+3(2048)(.177)+3(256)(3)<6007.
\end{align*}
Reflection preserves the energy norm. Linearity and the convexity of the
norm cover all segments, the two exact attachments, and the final
contraction to zero. The original two candidates coincide with
$\phi_{12}$ beyond minimum-endpoint radius $4096$, where its increments
are smaller than $10^{-6}$. Adding the certified correction radii gives
the common bound $.018$ used in \cref{sec:green}.
The supplied payload checker reproduces the array bounds from the two
frozen archives; the stationary correction radii still depend on the
separate full-space certificate.
